\newpage
\chapter*{Apéndices}
\addcontentsline{toc}{chapter}{Apéndices}
\thispagestyle{empty}
\newpage

\chapter{Interfaz de ejecución de los benchmarks}
\label{anexo:interfaz}

Todos los benchmarks adaptados comparten la misma interfaz de línea de comandos,
lo que permite reutilizar la infraestructura de campaña sin cambios por
benchmark. La invocación general es:

\begin{verbatim}
python <BENCHMARK>_Python.py -c <CLASE> -t <HILOS> -m <MODO>
\end{verbatim}

donde \texttt{-c} selecciona la clase NAS (\texttt{S}, \texttt{W}, \texttt{A}),
\texttt{-t} fija el número de hilos del equipo OpenMP y \texttt{-m} selecciona el
modo de ejecución de OMP4Py (0~puro, 1~híbrido, 2~compilado, 3~compilado
tipado). La salida estándar incluye el tiempo NAS interno, los Mop/s y el estado
de verificación (\texttt{SUCCESSFUL} o \texttt{UNSUCCESSFUL}), que son los
valores que recopila la infraestructura de medición descrita en el
Capítulo~\ref{chap:metodologia}.


\chapter{Ejemplo de kernel adaptado: EP en modo tipado}
\label{anexo:kernel-ep}

El siguiente extracto corresponde al kernel principal de EP en su versión
tipada (modo~3), \texttt{ep\_compute\_types}. Ilustra las decisiones de
implementación descritas en la Sección~\ref{sec:decisiones-implementacion}: la
región \texttt{parallel} con reducción sobre \texttt{sx} y \texttt{sy}, la lista
explícita de variables privadas, las memoryviews tipadas
(\texttt{cython.double[:]}) y la sección crítica final que acumula el histograma.
Se ha abreviado el cuerpo aritmético del bucle interno por claridad.

\begin{verbatim}
@omp(compile=use_compiled_types())
def ep_compute_types(nn: int, nk: int, nk_plus: int, nq: int, an: float):
    sx: float = 0.0
    sy: float = 0.0
    ...
    q_values = numpy.repeat(0.0, nq)
    q: cython.double[:] = q_values

    with omp("parallel reduction(+:sx,sy) "
             "private(i,k,seed_iter,ik,kk,l,a1,x1_seed,t2_int,t4_int,"
             "t1,t3,a2,x2_seed,z,t1_local,t2_local,x1,x2,"
             "t1_value,t2_value,t3_value,t4_value)"):
        qq_values = numpy.repeat(0.0, nq)
        x_values = numpy.empty(nk_plus, dtype=numpy.float64)
        qq: cython.double[:] = qq_values
        x_local: cython.double[:] = x_values

        with omp("for schedule(static)"):
            for k in range(1, nn + 1):
                ...                       # generacion de muestras (Box-Muller)
                for i in range(nk):
                    x1 = 2.0 * x_local[2 * i] - 1.0
                    x2 = 2.0 * x_local[2 * i + 1] - 1.0
                    t1_value = x1 * x1 + x2 * x2
                    if t1_value <= 1.0:
                        t2_value = math.sqrt(-2.0 * math.log(t1_value)
                                             / t1_value)
                        t3_value = x1 * t2_value
                        t4_value = x2 * t2_value
                        l = int(max(abs(t3_value), abs(t4_value)))
                        qq[l] = qq[l] + 1.0
                        sx = sx + t3_value
                        sy = sy + t4_value

        with omp("critical"):
            for i in range(nq):
                q[i] = q[i] + qq[i]
\end{verbatim}

La declaración explícita de variables privadas en la cláusula
\texttt{private(...)} es la corrección descrita en la
Sección~\ref{sec:decisiones-implementacion}: sin ella, varias temporales se
compartían entre hilos y la verificación NAS fallaba.


\chapter{Infraestructura de ejecución en Finisterrae~III}
\label{anexo:infraestructura}

La campaña se automatiza con los scripts del directorio \texttt{scripts/ft3}. El
script \texttt{array\_run.sbatch} es la plantilla de Slurm que ejecuta cada
elemento del array de trabajos; cada elemento procesa varias filas de la matriz
de combinaciones de forma secuencial mediante \texttt{run\_one.py}.

\begin{verbatim}
#!/usr/bin/env bash
#SBATCH --job-name=omp4py-npb
#SBATCH --nodes=1
#SBATCH --ntasks=1
#SBATCH --cpus-per-task=32
#SBATCH --mem-per-cpu=4G
#SBATCH --time=06:00:00

source "$FT3_ROOT/scripts/ft3/env.sh"
ft3_export_python_runtime "${RUNNER_PYTHON_TAG:-3.15t}"
RUNNER_PYTHON="$(ft3_venv_python "${RUNNER_PYTHON_TAG:-3.15t}")"

start_index="$(( (SLURM_ARRAY_TASK_ID - 1) * rows_per_task + 1 ))"
end_index="$(( SLURM_ARRAY_TASK_ID * rows_per_task ))"
for index in $(seq "$start_index" "$end_index"); do
    "$RUNNER_PYTHON" "$FT3_ROOT/scripts/ft3/run_one.py" \
        --manifest "$MANIFEST" --index "$index" \
        --campaign-dir "$CAMPAIGN_DIR"
done
\end{verbatim}

La selección del intérprete \emph{free-threaded} y la desactivación explícita
del GIL se realizan en \texttt{env.sh} mediante la variable
\texttt{PYTHON\_GIL=0} y la ruta de bibliotecas correspondiente a cada versión:

\begin{verbatim}
ft3_export_python_runtime() {
    local prefix; prefix="$(ft3_python_prefix "$1")"
    export LD_LIBRARY_PATH="$prefix/lib${LD_LIBRARY_PATH:+:$LD_LIBRARY_PATH}"
    export PYTHON_GIL=0
}
\end{verbatim}


\chapter{Estructura de los resultados de campaña}
\label{anexo:resultados-estructura}

Cada campaña genera un directorio bajo \texttt{results/} con los siguientes
artefactos, que constituyen la base de las tablas y figuras de los
Capítulos~\ref{chap:resultados},~\ref{chap:profiling}
y~\ref{chap:versiones}:

\begin{itemize}
  \item \texttt{manifest.csv}: matriz completa de combinaciones de la campaña.
  \item \texttt{raw/}: registros crudos con nombre semántico (versión de Python,
        benchmark, clase, modo, número de hilos y repetición).
  \item \texttt{records/}: un fichero JSON por ejecución, con tiempo NAS interno,
        tiempo externo de proceso, Mop/s, código de retorno y verificación.
  \item \texttt{summary.csv}: tabla plana con todas las ejecuciones.
  \item \texttt{summary\_best.csv}: agregados de mejor y media por combinación.
  \item \texttt{env.txt}: metadatos de la campaña (versiones de Python, NumPy,
        Cython y variables de entorno).
\end{itemize}

El registro de la versión exacta de cada componente y de las variables de
entorno responde a la implicación metodológica del
Capítulo~\ref{chap:versiones}: en evaluaciones de Python \emph{free-threaded} la
configuración del runtime y la afinidad de hilos pueden alterar la conclusión
experimental, por lo que deben quedar documentadas junto con los resultados.
